\chapter{The mfact Manufacturing Framework}
\label{ch:mfact}
\label{sec:mfact}

Chapter~\ref{ch:background} situated this thesis against Lean's ecosystem
and the wider LLM-assisted-proving literature. This chapter turns to the
apparatus the thesis actually builds: \textsf{mfact}, the Lean~4 layer
that gives the manufacturing pipeline's epistemics --- candidate,
refusal, admission, manifest, certified release, valid objection --- a
first-class typed representation rather than leaving them as informal
convention. As paper/main.tex states it, \textsf{mfact} is deliberately
small: seven modules and a command-line interface. It is not the whole
system. It is the Lean/Lake standing apparatus that sits inside a larger
standing-manufacturing graph whose other rails (process evidence,
benchmark evidence, correspondence evidence, publication evidence) are
developed in later chapters. This chapter develops, in order, the
governing law the whole pipeline answers to; the typed vocabulary of
candidates and refusals; admission with a covers proof; the manifest,
certified release, and closed-objection theorem; the ledger law that
locates authority in a machine-checkable ledger rather than in file
paths or naming conventions; and the agent actuation constitution that
governs how any producer --- human or machine --- is permitted to touch
this apparatus.

\section{The Receipted Manufacturing Law}
\label{sec:mfact-receipted-law}

The governing law of the pipeline, stated once in
Chapter~\ref{ch:intro} and elaborated here, is
\[
A = \mu_B(O^*_B), \qquad R_B = \mathrm{receipt}(A).
\]
Every symbol names a distinct, checkable object rather than a stage of
prose. $B$ is the declared boundary: the Lean packages, the declaration
catalog, the manifest schema, and the audited theorem set that together
fix what counts as ``inside'' for a given release. $O$ is raw
observation --- candidate output from a human editor, a language model,
a ggen template, or a repository fragment, with no standing of its own.
$O^*_B$ is admitted observation inside $B$: typed, kernel-checked,
bounded, and replayable, the only class of observation the manufacturing
transform $\mu_B$ is permitted to consume. $\mu_B$ is the lawful
manufacturing transform itself --- ggen rendering, Lean kernel
admission, the Lake build, and the axiom/sorry audit, all parameterized
by the boundary $B$ they operate inside. $A$ is the resulting artifact
with standing: an admitted theorem, a module, a manifest entry, a
release. $R_B$ is the receipt --- the build log, the manifest hash, the
axiom audit output, the negative-fixture result, the replay evidence ---
that proves $A$ is a genuine consequence of $\mu_B$ acting on $O^*_B$,
rather than merely an assertion made about $A$ after the fact. The
asymmetry this equation encodes is the one already stated in
Chapter~\ref{ch:intro}: raw generated output is $O$, never $O^*$, until
it has passed through an admission step; manufacture produces the
artifact, the receipt proves the consequence relation held, and replay
proves the receipt was non-ornamental --- that a third party, not just
the original run, can recompute the same result. As paper/main.tex puts
it, the pipeline exists to refuse the gap between contact and
consequence.

Within this law the manufacturing graph distinguishes four kinds of
evidence, each answering to a different admission boundary rather than
being an interchangeable substitute for the others: \emph{formal}
evidence (Lean theorems and their axiom audits, the subject of this
chapter and Chapter~\ref{ch:procint}); \emph{operational} evidence (the
praxis-graphlaw benchmark and its execution validation via rslab,
Chapter~\ref{ch:empirical}); \emph{correspondence} evidence
(Aeneas-extracted implementation image related to admitted \textsf{procint}
semantics, Chapter~\ref{ch:correspondence}); and \emph{publication}
evidence (generated paper fragments and manifests, of which this
thesis's own numeric claims are themselves instances). The
\texttt{final\_status.tex} fragment below reports the current release's
standing across several of these evidence kinds at once, rendered
directly from the manifest rather than typed here.

\begin{table}[h]
\centering
\input{../paper/final_status}
\caption{Current release standing, rendered from the manifest via \texttt{final\_status.tex}.}
\label{tab:mfact-final-status}
\end{table}

In this architecture ggen is not a code generator in the ordinary sense.
It is the projection half of $\mu_B$: it maps a declared obligation
graph into candidate Lean surfaces. The Lean kernel supplies the
admission half. Projection without admission confers no standing;
admission without a projection mechanism does not scale to a
several-hundred-declaration corpus. Neither half is optional, and
neither half does the other's job --- a point returned to below when
Section~\ref{sec:mfact-ledger} distinguishes what a template may render
from what only kernel admission may certify.

\paragraph{Rice's boundary and manufactured decidability.} Rice's
theorem bars a nontrivial semantic decision procedure over arbitrary
programs, and the manufacturing law does not evade this boundary; it
changes the object of judgment. The pipeline does not decide the
semantic correctness of an arbitrary generated artifact considered in
isolation. It decides whether a bounded candidate has been admitted by a
specific Lean kernel, under a declared axiom policy, with a complete
manifest receipt --- a decidable question that replaces an undecidable
one rather than answering it. Undecidability is fenced, not defeated,
and the standing this apparatus manufactures is standing \emph{inside
that fence}, a scope this chapter and Chapter~\ref{ch:discussion} both
return to when discussing what the pipeline does not claim.

\paragraph{Status is algebra, not prose.} The release distinguishes a
closed vocabulary of standing values --- \textsc{Declared},
\textsc{Extracted}, \textsc{Stated}, \textsc{Proven}, \textsc{Refused},
\textsc{Blocked}, \textsc{BuildBroken}, \textsc{Unknown},
\textsc{Unsupported}, and \textsc{Alive} --- and no transition between
them is implicit or inferred from surrounding prose. \textsc{Extracted}
is not proof; \textsc{Unknown} is not admitted; \textsc{Unsupported} is
not the same refusal as \textsc{Refused}. A candidate artifact, whatever
its origin, remains exactly that --- a candidate --- until an explicit
admission step assigns it one of these values, and the repository's
broader status taxonomy (\textsc{Alive}, \textsc{Partial\_Alive},
\textsc{Blocked}, \textsc{Blocked\_External}, \textsc{Build\_Broken},
\textsc{Refused}, \textsc{Pending\_External\_Actuation}, \textsc{Valid},
\textsc{Unsupported}, \textsc{Planned}, \textsc{Stated},
\textsc{Replay\_Not\_Run}, \textsc{In\_Progress}, per AGENTS.md) refines
this same discipline at the level of whole packets and reports, not just
individual theorems.

\section{Candidates and the Typed Refusal Family}
\label{sec:mfact-candidates}

A \texttt{Candidate} in \textsf{mfact} is content plus provenance ---
producer identity and a content hash --- reflecting that every producer,
human or automated, is untrusted by construction with respect to the
formal corpus; nothing about a candidate's origin changes what it is
permitted to do. A failed admission is not represented as an untyped
error string but as a value of a closed inductive type,
\texttt{Refusal}, whose constructors --- \texttt{kernelRejected},
\texttt{sorryPresent}, \texttt{unauthorizedAxiom},
\texttt{missingEvidence}, \texttt{malformedModel},
\texttt{invalidFixture} --- name specific, checkable failure modes.
Because the type is closed and carries no default case, an unrecognized
failure is itself required to be classified as a refusal rather than
silently passed through; there is no code path by which a candidate can
fail to be admitted without producing a typed, loggable, re-queueable
record of why.

This same discipline is carried into the repository's engineering
practice at a coarser grain, as a vocabulary of typed refusals a human
or agent contributor invokes when a proposed change would confer
standing without a corresponding admission step:
\texttt{HAND\_CODED\_GENERATED\_OUTPUT} (a generated artifact was
hand-edited directly), \texttt{GENERATED\_OUTPUT\_DRIFT} (a rendered
file no longer matches its declared source), \texttt{MISSING\_GGEN\_SOURCE}
and \texttt{MISSING\_GGEN\_TEMPLATE} (an artifact is ledgered as
generated but its producing source or template cannot be found),
\texttt{ORPHAN\_GENERATED\_FILE} and \texttt{UNREGISTERED\_PAPER\_FRAGMENT}
(a file carries the shape of a generated or paper artifact but is absent
from the ledger), \texttt{UNSUPPORTED\_STANDING\_CLAIM} and
\texttt{STATED\_PROMOTED\_TO\_PROVEN} (a claim exceeds its evidence, or a
status transition happened without the checked guard that must license
it), \texttt{MANUAL\_RELEASE\_COUNT} and \texttt{MANUAL\_RELEASE\_HASH}
(a release number or hash was typed by hand instead of read from a
generated fragment), and \texttt{RECEIPT\_RECURSION\_REFUSED} (an
operation would mutate the identity of a release it is supposed to be
reporting on, discussed further in
Section~\ref{sec:mfact-core-identity}). A further code,
\texttt{SOURCE\_CHANGE\_ASSERTION\_UNSUPPORTED}, is planned to require
that any claim of a changed source ship with an actual changed
source, template, or TTL fragment in the same commit, not a bare
assertion that a change occurred. Section~\ref{sec:ai-guardrail-audit} of
Chapter~\ref{ch:ai-development} discusses two refusal codes specific to
one family of deliberately unfinished theorems ---
\texttt{COUNTERMODEL\_PROMOTION\_REFUSED} and its associated guard name,
\texttt{countermodel\_not\_promoted} --- as an instance of this same
family, added in direct response to an audited failure rather than
speculatively.

\section{Admission with a Covers Proof}
\label{sec:mfact-admission}

The move from candidate to artifact is not a boolean flag but a bundled
record. An \texttt{Admission} gathers the evidence for one candidate:
the kernel acceptance record, the sorry census, the axiom-audit output,
and a \texttt{covers} field --- a proof term witnessing that the
collected evidence entries jointly cover the obligations the admission
gate imposes. This last field is the structural device that prevents an
admission from being assembled with a silently missing check: the type
of \texttt{covers} forces every required obligation to be discharged by
some evidence entry before an \texttt{Admission} value can be
constructed at all, so ``we forgot to check the axiom footprint'' is not
a runtime bug this framework can express, only a compile-time
impossibility. The gate function, \texttt{admit}, is total: applied to a
candidate it returns either an \texttt{Admission} or a \texttt{Refusal},
with no third outcome such as a pending or ambiguous state. Totality
here is doing real epistemic work --- it is what makes ``the candidate
was neither admitted nor refused'' an unrepresentable claim rather than
merely a discouraged one.

\section{Manifest, Certified Release, and the Closed Objection Family}
\label{sec:mfact-manifest}

A \texttt{Manifest}, serializable to JSON, lists artifacts together with
their content hashes, axiom sets, fixture results, and the proven or
stated status of each declaration --- the object rendered, for this
release, into Table~\ref{tab:mfact-final-status} above and into the
per-declaration ladder discussed further in Chapter~\ref{ch:empirical}.
A \texttt{CertifiedRelease} pairs a manifest with its audit obligations.
Against a release \texttt{R}, \texttt{ValidObjection R} is a closed
inductive family whose constructors are the \emph{only} admissible forms
an objection to \texttt{R} may take --- an unaudited theorem, a missing
content hash, an unauthorized axiom, a failing fixture --- and each
constructor demands its own evidence to be instantiated, so an objection
cannot itself be raised without being checkable. The per-release theorem
\begin{center}
\texttt{theorem no\_valid\_objection : IsEmpty (ValidObjection R)}
\end{center}
is discharged by refuting each constructor of \texttt{ValidObjection}
against the corresponding checked field of \texttt{R}: there is no
unaudited theorem because the audit ran, no missing hash because the
manifest is complete, and so on for each constructor in turn. For the
release this thesis reports, the theorem is proven and, per
\texttt{\#print axioms}, depends on no axioms at all --- not even the
three standard Lean/Mathlib axioms (\texttt{propext},
\texttt{Classical.choice}, \texttt{Quot.sound}) that most of
\textsf{procint}'s own theorems carry. This is deliberately the
strongest possible case in the release: the framework's own
closed-objection theorem is proven from definitional unfolding alone.
The theorem's scope is exactly its constructors, and this thesis is
careful not to overstate it: \texttt{no\_valid\_objection} says the
\emph{enumerated} objection forms are refuted against this release's
checked fields, not that the release is beyond all possible criticism,
including criticism outside the closed family's four constructors (for
example, a citation-fidelity objection of the kind STANDING.md
explicitly flags as outside this pipeline's scope in its own Scope
section --- a limitation returned to in Chapter~\ref{ch:discussion}).

\subsection{The Boundary of Present Standing}
\label{sec:mfact-boundary}

Because status is algebra and not prose, this framework is also
positioned to state precisely, rather than gesture at, what remains
outside \textsc{Proven} in the current catalog. Of the current release's
\StatedCount{} stated declarations, two categories are worth naming
explicitly because they have opposite epistemic character. Five
declarations in the workflow countermodel family ---
\texttt{crownCounter\_reaches\_mid}, \texttt{crownCounter\_reaches\_final},
\texttt{crownCounter\_sound}, \texttt{crownCounter\_not\_bounded}, and
\texttt{WfNet.infinite\_transition\_countermodel\_sound\_not\_bounded} ---
are \emph{deliberately and permanently} stated: they witness why the
crown-jewel soundness equivalence (discussed in
Chapter~\ref{ch:procint}) requires the hypothesis \texttt{[Finite T]},
by exhibiting a workflow net with infinitely many transitions that is
sound but whose short-circuited net is unbounded. These are not
unfinished proof obligations queued for future work; they are a negative
witness, and Chapter~\ref{ch:ai-development} describes the mechanical
guard that refuses any attempt to promote them to proven standing. By
contrast, exactly one declaration in the current \TotalDecls-declaration
catalog is genuinely open proof work:
\texttt{ProcInt.BranchingProcess.isUnfoldingOf\_statement}, an
unfolding-correctness statement tracked separately in
\texttt{research/wfnet/obligations.toml} as obligation order~3
(\texttt{unfolding\_correctness}). Its statement, as recorded in
\texttt{packs/lean-math-pack/fragments/petriext.ttl}, attributes the
homomorphism condition it formalizes to Engelfriet's 1991 treatment of
branching-process unfoldings\footnote{The declaration's docstring reads,
verbatim: ``a branching process is a legal unfolding of net N when its
occurrence net is acyclic, has no backward conflict, and the labeling
respects the pre/post structure of N on singleton arcs (Engelfriet 1991,
Def.~3.1 homomorphism condition, specialized to ordinary nets).'' No
bibliographic entry for this attribution exists in \texttt{paper/refs.bib};
this thesis reports the attribution as the source states it, in a
footnote rather than as a \texttt{\textbackslash cite}, rather than
inventing a bibtex record to close the gap.}. This distinction --- a
deliberately closed negative witness versus a single, honestly reported
open item --- is itself an instance of the status algebra this section
opened with: neither is described here as a gap in the ordinary sense,
because only one of the two actually is one.

\section{The Ledger Law: Authority from the Manifest, Not the Path}
\label{sec:mfact-ledger}

A recurring temptation in a manufactured-artifact repository is to infer
a file's authority from where it lives or how it is named --- a
\texttt{generated/} directory, a header comment reading ``DO NOT
EDIT.'' This repository's ledger law rejects that inference explicitly.
There are, in its own words, no generated files; there are only
artifacts with receipts, and every repository file is first-class at its
canonical path. Authority comes from the artifact ledger,
\texttt{.mfact/artifacts.toml}, and from nowhere else. Concretely, this
has two directions of consequence. If a file is ledgered as produced by
ggen or a builder script, it may not be patched directly as a final
solution to a problem it exhibits: the sources or templates that produce
it must be modified, the file re-rendered, and \texttt{just regen-check}
must pass again before the edit is considered done --- an unreplayable
hand-edit to a ledgered file is \texttt{ARTIFACT\_DRIFT\_REFUSED}
regardless of how small or how obviously correct the hand-edit looked.
If, conversely, a file is \emph{not} ledgered but nonetheless carries
release standing, counts, audit status, or certification data, it is
classified as \texttt{ORPHAN\_ARTIFACT\_REFUSED}: the task is to either
ledger it properly or refuse to proceed, never to let an unregistered
file quietly carry standing.

Two surfaces are explicitly and permanently excluded from this ledger,
by design rather than by omission: \texttt{procint/Playground/**}, a
hand-authored Lean demonstration surface that is never ggen-rendered and
carries no standing claims of its own, and \texttt{pylab/**}, the
analogous hand-authored Python research surface. Both are developed at
length in Chapter~\ref{ch:ai-development}, which also describes a
mechanized scanner, itself living in \texttt{pylab/}, whose job is to
detect exactly the two ledger-law violations named in this section ---
drifted ledgered artifacts and unledgered orphans --- as a read-only,
repeatable check rather than a one-time manual audit.

\section{The Agent Actuation Constitution}
\label{sec:mfact-actuation}

Because candidates in this repository include agent-produced patches at
every layer --- TTL fragments, ggen templates, builder scripts, and this
thesis's own prose --- the repository constrains not just what may be
edited but how any change is permitted to take effect. Agents actuate
only through \texttt{just} recipes; raw invocations of \texttt{lake},
\texttt{ggen}, \texttt{mfact}, or git as a final actuation path are
disallowed unless a recipe or a doctor report explicitly instructs
otherwise for debugging. If a task requires an actuation path that does
not yet exist as a recipe, the constitution requires adding the recipe
first and only then using it, rather than reaching for the underlying
tool directly. This is not merely a style preference: a raw command
invoked outside its recipe bypasses whatever bookkeeping the recipe
performs around it (writing an ephemeral report, checking a
precondition, chaining a dependent step), reintroducing exactly the kind
of untracked side effect the ledger law in
Section~\ref{sec:mfact-ledger} exists to prevent.

The recipe surface itself is split into two classes with different
mutation rights. Diagnostic recipes --- \texttt{status}, \texttt{next},
\texttt{doctor}, \texttt{trace}, \texttt{why}, \texttt{theorem-status},
\texttt{proof-blockers}, \texttt{fixtures}, \texttt{docs-check} --- are
read-only: none of them may dirty the working tree, and none of their
output is itself a certified artifact. Only \texttt{report-write},
\texttt{check}, and \texttt{release} write the ephemeral cockpit reports
under \texttt{.mfact/reports/latest.*}, and these reports are themselves
gitignored and never ledgered --- they are a working surface for the
agent driving the recipe, not a standing claim. Certified status, by
contrast, lives only in the \texttt{release/} artifacts produced by the
manufacturing recipes proper: \texttt{render}, \texttt{build},
\texttt{audit}, \texttt{manifest}, \texttt{certify}, and
\texttt{standing-quadrature}. A completed task's report is required to
name the specific recipes it ran, not an ad hoc shell history, so that
the report itself is reconstructible from the recipe names alone.

\section{Core Release Identity and Recursion Refusal}
\label{sec:mfact-core-identity}

The manufacturing law of Section~\ref{sec:mfact-receipted-law} applies
recursively: the paper and this thesis are themselves artifacts of the
pipeline they describe, and their own claims must be receipted rather
than asserted. This creates a specific hazard the framework must guard
against explicitly. The core v26.7.7 release is frozen at tag
\texttt{\ReleaseTag}, with a manifest fold hash, a proven-declaration
count, and a total-declaration count fixed at that tagged commit ---
values this thesis reports throughout via the \texttt{release\_macros}
fragment rather than by typing them by hand. Post-release work,
including work on this thesis and on the correspondence and evaluation
material of later chapters, must not mutate the identity of the core
certified release it reports on. If an operation would change the core
fold hash, the proven count, the declaration count, or the manifest
itself, the constitution requires one of two outcomes and forbids the
third: either refuse the operation as
\texttt{RECEIPT\_RECURSION\_REFUSED}, or promote it explicitly into a
new core certification cycle with new manifest values recorded under a
new identity. What is not permitted is a silent third option in which
the core identity shifts underneath an unchanged tag. Post-release
witnesses such as \texttt{ProcInt.Release.PostRelease} are counted
separately, as \texttt{POST\_RELEASE\_WITNESSES}, and are never folded
into the core proven count; the core tag itself is pinned to the core
release commit and is not moved by later packet work, which receives its
own tags and its own packet hashes. \texttt{just certify} and
\texttt{just release} are required to verify that the release commit is
an ancestor of, or equal to, the certified tag before reporting a
passing \texttt{CERTIFIED\_RELEASE} status; a tag that predates the
currently rendered artifacts is treated as
\texttt{RECEIPT\_RECURSION\_REFUSED}, not filed away as a historical
curiosity to be noted and passed over. Chapter~\ref{ch:ai-development}
describes the mechanized check that implements this ancestry
verification and traces it to the specific 2026-07-07 audit finding that
motivated adding it.

\section{Summary}

This chapter has developed \textsf{mfact} as the typed apparatus that
makes the pipeline's receipted manufacturing law operational: candidates
and refusals as closed data types with no silent failure path;
admission as a total function whose \texttt{covers} proof makes an
incomplete check unrepresentable; the manifest, certified release, and
closed \texttt{ValidObjection} family as the objects a
\texttt{no\_valid\_objection} theorem is proven against, scoped exactly
to its enumerated constructors; the ledger law that locates authority in
\texttt{.mfact/artifacts.toml} rather than in file paths, headers, or
naming conventions; the agent actuation constitution that confines any
producer's actuation to named, classified \texttt{just} recipes; and the
core release identity's protection against self-referential mutation.
Chapter~\ref{ch:procint} now turns to the first substantial artifact
this apparatus was used to manufacture: \textsf{procint}, and in
particular the workflow-net soundness development whose crown-jewel
theorem and deliberate countermodel were introduced only by name in
Section~\ref{sec:mfact-boundary} above.
